\documentclass[12pt,a4paper]{ctexart}

\usepackage[margin=1in]{geometry}
\usepackage{amsmath,amssymb,amsthm,bm,mathtools}
\usepackage{hyperref}
\usepackage{enumitem}
\usepackage{booktabs}

\hypersetup{
    colorlinks=true,
    linkcolor=blue,
    citecolor=blue,
    urlcolor=blue
}

\newtheorem{theorem}{定理}[section]
\newtheorem{proposition}[theorem]{命题}
\newtheorem{assumption}[theorem]{假设}
\newtheorem{remark}[theorem]{说明}
\newtheorem{definition}[theorem]{定义}
\newtheorem{lemma}[theorem]{引理}
\newtheorem{corollary}[theorem]{推论}

\DeclareMathOperator{\col}{col}
\DeclareMathOperator{\diag}{diag}
\DeclareMathOperator{\atanh}{atanh}

\title{通信感知团队价值函数：\\
从组级博弈到结构化分解的完整推导}
\author{}
\date{}

\begin{document}

\maketitle
\tableofcontents
\newpage

% ============================================================
\section{问题建模}
% ============================================================

\subsection{单体动力学}

考虑 $n$ 追一的追逃博弈。每个智能体的状态为
\[
x = [p_1, p_2, p_3, v_1, v_2, v_3]^T \in \mathbb{R}^6,
\]
其中 $p = [p_1, p_2, p_3]^T$ 为三维位置，$v = [v_1, v_2, v_3]^T$ 为三维速度。
控制输入为
\[
u = [u_1, u_2, u_3]^T \in \mathbb{R}^3.
\]

动力学写为
\begin{equation}\label{eq:dynamics}
\dot{x} = f(x) + g(x)\,u,
\end{equation}
其中 $f\colon \mathbb{R}^6 \to \mathbb{R}^6$ 为漂移项，$g\colon \mathbb{R}^6 \to \mathbb{R}^{6 \times 3}$ 为输入增益矩阵。

\begin{assumption}[输入增益结构]\label{ass:g-structure}
输入增益矩阵的前三行为零：
\[
g(x) = \begin{bmatrix} \bm{0}_{3\times 3} \\ G(x) \end{bmatrix},
\qquad G(x) \in \mathbb{R}^{3\times 3},
\]
其中 $G(x)$ 在所考察区域内有界且非退化。此结构意味着控制输入直接作用于速度分量，不直接改变位置。
\end{assumption}

第 $j$ 个追逐者的动力学为
\begin{equation}
\dot{x}_j^p = f_j(x_j^p) + g_j(x_j^p)\,u_j^p, \qquad j = 1, \dots, n.
\end{equation}

逃避者的动力学为
\begin{equation}
\dot{x}^e = f_e(x^e) + g_e(x^e)\,u^e.
\end{equation}

\subsection{追踪误差与编队位移}

定义第 $j$ 个追逐者的追踪误差
\begin{equation}\label{eq:error}
\tilde{x}_j = x_j^p - x^e + r_j,
\end{equation}
其中 $r_j \in \mathbb{R}^6$ 为预设编队位移常量，定义追逐者 $j$ 相对于逃避者的期望最终偏移。记 $\tilde{p}_j$ 和 $\tilde{v}_j$ 分别为位置误差和速度误差分量。

追踪误差的动力学为
\begin{equation}\label{eq:error-dynamics}
\dot{\tilde{x}}_j = f_j(x_j^p) - f_e(x^e) + g_j(x_j^p)\,u_j^p - g_e(x^e)\,u^e.
\end{equation}

\subsection{通信模型}\label{sec:comm-model}

每个追逐者通过通信链路接收队友信息。定义追逐者 $k$ 在时刻 $t$ 广播的完整信息为
\begin{equation}\label{eq:comm-content}
\xi_k(t) = \bigl(x_k^p(t),\; u_k^p(t)\bigr) \in \mathbb{R}^9,
\end{equation}
包含状态（6 维）和当前控制输入（3 维）。

\begin{definition}[通信矩阵]\label{def:comm-matrix}
$\Gamma(t) = [\gamma_{kb}(t)] \in \{0,1\}^{n \times n}$，其中 $\gamma_{kb}(t) = 1$ 表示追逐者 $k$ 的信息在时刻 $t$ 能被追逐者 $b$ 接收。约定 $\gamma_{jj} = 1$。
\end{definition}

\begin{definition}[三级通信状态]\label{def:three-level}
对追逐者 $j$ 与队友 $k$，定义三级通信状态：
\begin{enumerate}[label=(\roman*)]
\item \textbf{全通信}：$j$ 同时收到 $x_k^p$ 和 $u_k^p$；
\item \textbf{状态通信}：$j$ 仅收到 $x_k^p$，$u_k^p$ 不可用；
\item \textbf{完全断联}：$j$ 未收到 $k$ 的任何信息。
\end{enumerate}
\end{definition}

\begin{remark}[通信控制输入的意义]
在离散时间实现中，$u_k^p(t)$ 为队友 $k$ 在时刻 $t$ 的加速度指令，该指令在 $t+\Delta t$ 之后才反映到 $x_k^p$ 的速度中。因此通信 $u_k^p$ 等价于提供一步预测：
$x_k^p(t+\Delta t) \approx x_k^p(t) + (f_k + g_k\,u_k^p)\,\Delta t$。
知道队友正在加速朝自己飞来还是远离，对闪避决策至关重要。
\end{remark}


% ============================================================
\section{组级博弈框架（理论动机）}
% ============================================================

\subsection{组级状态与性能指标}

将 $n$ 个误差块堆叠为组级状态
\[
Z = \col(\tilde{x}_1, \tilde{x}_2, \dots, \tilde{x}_n) \in \mathbb{R}^{6n},
\]
将追逐者控制堆叠为
\[
U^p = \col(u_1^p, u_2^p, \dots, u_n^p) \in \mathbb{R}^{3n}.
\]

定义运行代价
\begin{equation}\label{eq:group-cost}
\ell(Z, U^p, u^e) = \sum_{j=1}^{n} P_j(\tilde{x}_j) + \sum_{j=1}^{n} \Psi_{\bar{u}_p, R_1}(u_j^p) - \Psi_{\bar{u}_e, R_2}(u^e),
\end{equation}
其中状态代价为
\begin{equation}
P_j(\tilde{x}_j) = \left(\frac{\tilde{x}_j}{s}\right)^{\!T} Q \left(\frac{\tilde{x}_j}{s}\right), \qquad Q \succ 0,
\end{equation}
$s \in \mathbb{R}^6$ 为状态缩放向量。非二次输入代价为
\begin{equation}\label{eq:nonquad}
\Psi_{\bar{u}, R}(u) = \sum_{\ell=1}^{3} 2\bar{u}\,r_\ell \left[ u_\ell\,\atanh\!\left(\frac{u_\ell}{\bar{u}}\right) + \frac{\bar{u}}{2} \ln\!\left(1 - \frac{u_\ell^2}{\bar{u}^2}\right) \right].
\end{equation}

组级性能指标为
\[
J(Z(0), U^p, u^e) = \int_0^\infty \ell(Z, U^p, u^e)\,dt.
\]

\subsection{组级 HJI 方程与最优控制律}

组级值函数
$V^*(Z) = \min_{U^p} \max_{u^e} J$
满足 HJI 方程
\begin{equation}\label{eq:group-hji}
0 = \min_{U^p} \max_{u^e} \left[\ell(Z, U^p, u^e) + \nabla V^*(Z)^T \dot{Z}\right].
\end{equation}

由鞍点条件可推导（详见附录~\ref{app:control}）：

\textbf{追逐者}（第 $j$ 个）：
\begin{equation}\label{eq:group-up}
u_j^{p*} = -\bar{u}_p\,\tanh\!\left(\frac{1}{2\bar{u}_p}\,R_1^{-1}\,g_j^T \frac{\partial V^*}{\partial \tilde{x}_j}\right).
\end{equation}

\textbf{逃避者}：
\begin{equation}\label{eq:group-ue}
u^{e*} = -\bar{u}_e\,\tanh\!\left(\frac{1}{2\bar{u}_e}\,R_2^{-1}\,g_e^T \sum_{j=1}^{n} \frac{\partial V^*}{\partial \tilde{x}_j}\right).
\end{equation}

\subsection{计算不可行性——维度灾难}

\begin{remark}[核心困难]\label{rmk:curse}
$V^*(Z)$ 定义在 $\mathbb{R}^{6n}$ 上。对 $n = 4$ 为 24 维，$n = 6$ 为 36 维。若采用近似
$\hat{V}(Z) = \hat{W}^T \phi(Z)$，
基函数 $\phi(Z)$ 必须在 $\mathbb{R}^{6n}$ 上有足够表达力。原文用于 $\mathbb{R}^6$ 的 6 项基函数无法直接扩展至高维。

更关键的是：若 $\phi(Z)$ 采用块对角结构 $\phi(Z) = \col(\phi(\tilde{x}_1), \dots, \phi(\tilde{x}_n))$，则 $\partial\hat{V}/\partial\tilde{x}_j$ 仅取决于 $\tilde{x}_j$ 本身，退化为独立 critic，丧失一切协调能力。
\end{remark}


% ============================================================
\section{结构化近似：通信感知团队价值函数}
% ============================================================

\subsection{结构化分解的思路}

为在保留理论性质的同时回避维度灾难，我们将 $V^*(Z)$ 近似分解为
\begin{equation}\label{eq:total-value}
\hat{V}(Z) = \sum_{j=1}^{n} V_j^{\mathrm{team}},
\end{equation}
其中每个追逐者的团队价值函数为
\begin{equation}\label{eq:team-value}
\boxed{
V_j^{\mathrm{team}}(\tilde{x}_j,\, \{\xi_k\}_{k \neq j})
= W_j^T \phi(\tilde{x}_j) + W_c^T \psi_j(x_j^p,\, \{\xi_k\}_{k \neq j})
}
\end{equation}
其中：
\begin{itemize}
\item $W_j \in \mathbb{R}^{n_\phi}$：第 $j$ 个追逐者的\textbf{个体权重}（每人独立）；
\item $W_c \in \mathbb{R}^{n_c}$：\textbf{耦合权重}（全队共享）；
\item $\phi(\tilde{x}_j) \in \mathbb{R}^{n_\phi}$：个体特征，仅依赖自身误差；
\item $\psi_j \in \mathbb{R}^{n_c}$：耦合特征，依赖自身状态与通信接收的队友状态和控制 $\xi_k = (x_k^p, u_k^p)$。
\end{itemize}

\begin{remark}[为什么 $W_c$ 全队共享]
耦合特征 $\psi_j$ 已经编码了追逐者特定信息（不同的队友集合和相对位置），因此 $W_c$ 共享不损失个性。共享的优势：(i) 参数量从 $9n$ 降至 $9$；(ii) 所有追逐者样本汇入同一 LS 问题，条件数更优；(iii) 对称协调——无追逐者有特殊地位。
\end{remark}

\subsection{个体特征 $\phi(\tilde{x}_j)$}

与原文一致，采用 $n_\phi = 6$ 项耦合基函数。记缩放向量 $\sigma \in \mathbb{R}^3_{>0}$，增益 $\alpha > 0$：
\begin{equation}\label{eq:phi}
\phi(\tilde{x}_j) = \alpha \begin{bmatrix}
\tilde{p}_{j,1}\,\tilde{v}_{j,1} / \sigma_1 \\
\tilde{p}_{j,2}\,\tilde{v}_{j,2} / \sigma_2 \\
\tilde{p}_{j,3}\,\tilde{v}_{j,3} / \sigma_3 \\
\tfrac{1}{2}\,\tilde{v}_{j,1}^2 \\
\tfrac{1}{2}\,\tilde{v}_{j,2}^2 \\
\tfrac{1}{2}\,\tilde{v}_{j,3}^2
\end{bmatrix} \in \mathbb{R}^6.
\end{equation}

其 Jacobian $J_\phi = \partial\phi/\partial\tilde{x}_j \in \mathbb{R}^{6 \times 6}$：
\begin{equation}\label{eq:jac-phi}
J_\phi = \alpha \begin{bmatrix}
\tilde{v}_{j,1}/\sigma_1 & 0 & 0 & \tilde{p}_{j,1}/\sigma_1 & 0 & 0 \\
0 & \tilde{v}_{j,2}/\sigma_2 & 0 & 0 & \tilde{p}_{j,2}/\sigma_2 & 0 \\
0 & 0 & \tilde{v}_{j,3}/\sigma_3 & 0 & 0 & \tilde{p}_{j,3}/\sigma_3 \\
0 & 0 & 0 & \tilde{v}_{j,1} & 0 & 0 \\
0 & 0 & 0 & 0 & \tilde{v}_{j,2} & 0 \\
0 & 0 & 0 & 0 & 0 & \tilde{v}_{j,3}
\end{bmatrix}.
\end{equation}


\subsection{耦合特征 $\psi_j$ 的设计}\label{sec:coupling-feature}

定义追逐者 $j$ 与队友 $k$ 之间的相对状态
\[
\Delta x_{jk} = x_j^p - x_k^p = [\Delta p_{jk}^T,\; \Delta v_{jk}^T]^T \in \mathbb{R}^6.
\]

定义\textbf{单对耦合特征映射} $\phi_c\colon \mathbb{R}^6 \times \mathbb{R}^3 \to \mathbb{R}^{n_c}$：
\begin{equation}\label{eq:phi-c}
\phi_c(\Delta x,\, u_k) = \begin{bmatrix} \phi_s(\Delta x) \\ \phi_u(\Delta v, u_k) \end{bmatrix},
\end{equation}
由两部分组成。

\paragraph{状态耦合 $\phi_s \in \mathbb{R}^6$}
编码追逐者间的相对运动趋势：
\begin{equation}\label{eq:phi-s}
\phi_s(\Delta x) = \gamma_s \begin{bmatrix}
\Delta p_1\,\Delta v_1 / \sigma_1^c \\
\Delta p_2\,\Delta v_2 / \sigma_2^c \\
\Delta p_3\,\Delta v_3 / \sigma_3^c \\
\tfrac{1}{2}\,\Delta v_1^2 \\
\tfrac{1}{2}\,\Delta v_2^2 \\
\tfrac{1}{2}\,\Delta v_3^2
\end{bmatrix},
\end{equation}
其中 $\gamma_s > 0$ 为状态耦合增益，$\sigma^c \in \mathbb{R}^3_{>0}$ 为追逐者间位置缩放。

物理含义：$\Delta p_\ell \cdot \Delta v_\ell > 0$ 表示两追逐者在第 $\ell$ 轴上正在分离；$< 0$ 表示正在接近。$\Delta v_\ell^2$ 编码相对速度大小。

\paragraph{控制耦合 $\phi_u \in \mathbb{R}^3$}
编码队友控制意图与相对速度的交互：
\begin{equation}\label{eq:phi-u}
\phi_u(\Delta v, u_k) = \gamma_u \begin{bmatrix}
\Delta v_1 \cdot u_{k,1} \\
\Delta v_2 \cdot u_{k,2} \\
\Delta v_3 \cdot u_{k,3}
\end{bmatrix},
\end{equation}
其中 $\gamma_u > 0$ 为控制耦合增益。

\begin{remark}[为什么采用 $\Delta v \cdot u_k$ 而非 $\Delta p \cdot u_k$]\label{rmk:why-dv-uk}
由假设~\ref{ass:g-structure}，$g(x)$ 前三行为零，故控制律仅取决于 $\partial V_j/\partial x_j$ 的\textbf{速度分量}（第 4--6 个分量）。

特征 $\Delta v_\ell \cdot u_{k,\ell}$ 对 $\Delta v_\ell$（即 $v_{j,\ell}$，因 $v_{k,\ell}$ 为通信接收常量）的偏导为 $\gamma_u u_{k,\ell}$，位于速度分量，可通过 $g^T$ 直接进入控制律。

而 $\Delta p_\ell \cdot u_{k,\ell}$ 的梯度位于\textbf{位置分量}（第 1--3 个），被 $g^T$ 乘以零，无法进入控制律。
\end{remark}

总耦合特征维度为 $n_c = 6 + 3 = 9$。

\paragraph{耦合特征的聚合}
取可通信队友的平均：
\begin{equation}\label{eq:psi-j}
\psi_j = \frac{1}{|\mathcal{N}_j|} \sum_{k \in \mathcal{N}_j} \phi_c(\Delta x_{jk},\, u_k^p),
\end{equation}
其中 $\mathcal{N}_j = \{k : k \neq j,\;\gamma_{kj}(t) = 1\}$。当 $\mathcal{N}_j = \varnothing$ 时，$\psi_j = \bm{0}$。

\begin{remark}[三级通信下的耦合特征]
\begin{enumerate}[label=(\roman*)]
\item 全通信：$\phi_c = [\phi_s^T,\, \phi_u^T]^T$，$n_c = 9$ 维完整激活；
\item 状态通信：$u_k$ 不可用，$\phi_u = \bm{0}$，等效 6 维；
\item 完全断联：$\psi_j = \bm{0}$，耦合项消失。
\end{enumerate}
\end{remark}


% ============================================================
\section{最优控制律推导}
% ============================================================

\subsection{团队价值函数的梯度}

控制律需要 $\partial V_j^{\mathrm{team}}/\partial x_j^p$。由 $\tilde{x}_j = x_j^p - x^e + r_j$ 得 $\partial\tilde{x}_j/\partial x_j^p = I_6$，故

\begin{equation}\label{eq:grad-V-team}
\boxed{
\frac{\partial V_j^{\mathrm{team}}}{\partial x_j^p}
= \underbrace{\left(\frac{\partial \phi}{\partial \tilde{x}_j}\right)^{\!T} W_j}_{\text{个体梯度}}
+ \underbrace{\left(\frac{\partial \psi_j}{\partial x_j^p}\right)^{\!T} W_c}_{\text{耦合梯度}}
}
\end{equation}

在控制律的计算中，队友信息 $\xi_k = (x_k^p, u_k^p)$ 为通信接收的\textbf{给定量}，不对其求导。

\subsection{耦合特征的 Jacobian}\label{sec:jac-coupling}

由 \eqref{eq:psi-j}，
\[
\frac{\partial \psi_j}{\partial x_j^p}
= \frac{1}{|\mathcal{N}_j|} \sum_{k \in \mathcal{N}_j} J_c(\Delta x_{jk},\, u_k^p) \in \mathbb{R}^{n_c \times 6},
\]
其中
\begin{equation}\label{eq:Jc}
J_c = \frac{\partial \phi_c}{\partial(\Delta x)} = \begin{bmatrix} J_s \\ J_u \end{bmatrix}.
\end{equation}

\textbf{状态耦合 Jacobian} $J_s \in \mathbb{R}^{6 \times 6}$：
\begin{equation}\label{eq:Js}
J_s = \gamma_s \begin{bmatrix}
\Delta v_1/\sigma_1^c & 0 & 0 & \Delta p_1/\sigma_1^c & 0 & 0 \\
0 & \Delta v_2/\sigma_2^c & 0 & 0 & \Delta p_2/\sigma_2^c & 0 \\
0 & 0 & \Delta v_3/\sigma_3^c & 0 & 0 & \Delta p_3/\sigma_3^c \\
0 & 0 & 0 & \Delta v_1 & 0 & 0 \\
0 & 0 & 0 & 0 & \Delta v_2 & 0 \\
0 & 0 & 0 & 0 & 0 & \Delta v_3
\end{bmatrix}.
\end{equation}

\textbf{控制耦合 Jacobian} $J_u \in \mathbb{R}^{3 \times 6}$：
\begin{equation}\label{eq:Ju}
J_u = \gamma_u \begin{bmatrix}
0 & 0 & 0 & u_{k,1} & 0 & 0 \\
0 & 0 & 0 & 0 & u_{k,2} & 0 \\
0 & 0 & 0 & 0 & 0 & u_{k,3}
\end{bmatrix}.
\end{equation}

$J_u$ 的非零项全部位于第 4--6 列（速度分量），因此经 $g^T$ 提取后确实进入控制律。

\begin{remark}[控制耦合进入控制律的完整路径]\label{rmk:control-path}
记 $W_c = [W_c^{s,T},\, W_c^{u,T}]^T$，其中 $W_c^s \in \mathbb{R}^6$，$W_c^u \in \mathbb{R}^3$。耦合梯度的速度分量为
\begin{align*}
\left[\left(\frac{\partial \psi_j}{\partial x_j^p}\right)^{\!T} W_c\right]_{4:6}
&= \frac{1}{|\mathcal{N}_j|} \sum_{k \in \mathcal{N}_j} \Bigl[
\gamma_s\bigl(\Delta p_\ell / \sigma_\ell^c\bigr) W_{c,\ell}^s
+ \gamma_s\,\Delta v_\ell\, W_{c,\ell+3}^s \\
&\qquad\qquad\qquad\quad + \gamma_u\, u_{k,\ell}\, W_{c,\ell}^u
\Bigr]_{\ell=1,2,3}.
\end{align*}
最后一项 $\gamma_u\, u_{k,\ell}\, W_{c,\ell}^u$ 即为队友控制对自身控制律的影响：当队友在第 $\ell$ 方向加速（$u_{k,\ell} \neq 0$）且权重 $W_{c,\ell}^u$ 表明该方向需要协调时，追逐者 $j$ 的控制律做出响应。
\end{remark}

\subsection{追逐者控制律}

对第 $j$ 个追逐者，Hamiltonian 中与 $u_j^p$ 相关的驻值条件（推导见附录~\ref{app:control}）给出

\begin{equation}\label{eq:team-up}
\boxed{
u_j^{p*} = -\bar{u}_p\,\tanh\!\left(\frac{1}{2\bar{u}_p}\,R_1^{-1}\,g_j(x_j^p)^T\,
\frac{\partial V_j^{\mathrm{team}}}{\partial x_j^p}\right)
}
\end{equation}

将 \eqref{eq:grad-V-team} 代入展开：
\begin{equation}\label{eq:team-up-expanded}
u_j^{p*} = -\bar{u}_p\,\tanh\!\left(\frac{1}{2\bar{u}_p}\,R_1^{-1}\,g_j^T \left[
\left(\frac{\partial \phi}{\partial \tilde{x}_j}\right)^{\!T} W_j
+ \left(\frac{\partial \psi_j}{\partial x_j^p}\right)^{\!T} W_c
\right]\right).
\end{equation}

\subsection{逃避者控制律}

逃避者最大化追逐者团队的总代价。由 $\hat{V}(Z) = \sum_j V_j^{\mathrm{team}}$ 的组级 HJI 鞍点条件：
\begin{equation}\label{eq:team-ue}
u^{e*} = -\bar{u}_e\,\tanh\!\left(\frac{1}{2\bar{u}_e}\,R_2^{-1}\,g_e(x^e)^T \sum_{j=1}^{n} \frac{\partial V_j^{\mathrm{team}}}{\partial \tilde{x}_j}\right).
\end{equation}

\begin{remark}
追逐者控制 \eqref{eq:team-up} 仅使用\textbf{自身}梯度，实现分布式执行；逃避者控制 \eqref{eq:team-ue} 使用\textbf{全体}梯度之和，与组级 HJI~\eqref{eq:group-ue} 结构一致。
\end{remark}


% ============================================================
\section{代价函数设计}
% ============================================================

\subsection{个体运行代价}

第 $j$ 个追逐者的个体运行代价为
\begin{equation}\label{eq:cost-ind}
L_j^{\mathrm{ind}} = P_j(\tilde{x}_j) + \Psi_{\bar{u}_p, R_1}(u_j^p) - \frac{1}{n}\,\Psi_{\bar{u}_e, R_2}(u^e).
\end{equation}

注意逃避者代价在 $n$ 个追逐者间均分，使得 $\sum_j L_j^{\mathrm{ind}}$ 与组级代价 \eqref{eq:group-cost} 一致。

\subsection{分离惩罚函数}

\begin{remark}[为什么需要显式惩罚]
组级代价 \eqref{eq:group-cost} 中每一项 $P_j(\tilde{x}_j)$ 仅依赖第 $j$ 个追逐者自身的误差，追逐者之间无交叉项。即使两个追逐者走完全相同的路径，只要各自的 $\tilde{x}_j$ 趋零，代价同样低。优化器看不到路径重合。
\end{remark}

为显式激励路径分散，引入高斯核分离惩罚：
\begin{equation}\label{eq:sep-penalty}
S_j = \sum_{k \neq j} \exp\!\left(-\frac{\|p_j^p - p_k^p\|^2}{2\,d_{\mathrm{ref}}^2}\right),
\end{equation}
其中 $d_{\mathrm{ref}} > 0$ 控制惩罚作用范围。当追逐者 $j$ 与 $k$ 接近时，$S_j$ 增大。

\subsection{团队运行代价}

\begin{equation}\label{eq:cost-team}
L_j^{\mathrm{team}} = L_j^{\mathrm{ind}} + \gamma\,S_j,
\end{equation}
其中 $\gamma > 0$ 为分离惩罚权重。增量部分为
\begin{equation}\label{eq:cost-increment}
\Delta L_j \triangleq L_j^{\mathrm{team}} - L_j^{\mathrm{ind}} = \gamma\,S_j \ge 0.
\end{equation}


% ============================================================
\section{两阶段训练}
% ============================================================

\subsection{第一阶段：个体权重 $W_j$}

固定 $W_c = \bm{0}$（无耦合），训练每个追逐者的个体权重 $W_j$。此阶段与原文 V-SNAC 完全一致。

沿轨迹的 Bellman 差分方程：
\begin{equation}\label{eq:bellman-1}
\bigl(\phi(\tilde{x}_j(t+\Delta t)) - \phi(\tilde{x}_j(t))\bigr)^T W_j = -L_j^{\mathrm{ind}}(t) \cdot \Delta t.
\end{equation}

收集多步样本，通过岭回归最小二乘求解 $W_j \in \mathbb{R}^6$。

\begin{assumption}[第一阶段收敛]\label{ass:stage1-converge}
第一阶段训练收敛至 $W_j^*$。此假设与原文一致。
\end{assumption}

\subsection{第二阶段：耦合权重 $W_c$}

固定 $W_j = W_j^*$，训练共享耦合权重 $W_c$。

\paragraph{Bellman 差分方程的推导}
团队价值函数沿轨迹的时间导数为
\[
\frac{dV_j^{\mathrm{team}}}{dt} = \frac{d(W_j^T\phi)}{dt} + \frac{d(W_c^T\psi_j)}{dt}.
\]

在团队最优策略下 $dV_j^{\mathrm{team}}/dt = -L_j^{\mathrm{team}}$，而第一阶段保证 $d(W_j^T\phi)/dt \approx -L_j^{\mathrm{ind}}$，故
\[
\frac{d(W_c^T\psi_j)}{dt} \approx -L_j^{\mathrm{team}} + L_j^{\mathrm{ind}} = -\gamma\,S_j.
\]

离散化得

\begin{equation}\label{eq:bellman-2}
\boxed{
\bigl(\psi_j(t+\Delta t) - \psi_j(t)\bigr)^T W_c = -\gamma\,S_j(t) \cdot \Delta t
}
\end{equation}

\begin{remark}[样本量]
$W_c$ 全队共享，$n$ 个追逐者各提供 $T$ 步样本，LS 矩阵为 $nT \times 9$。例如 $n = 4$、$T = 250$，则有 1000 个样本训练 9 个参数。
\end{remark}

\begin{remark}[策略迭代]
第二阶段使用团队策略 \eqref{eq:team-up-expanded}（含当前 $W_c$）采集轨迹。$W_c$ 初始化为 $\bm{0}$，每次策略迭代更新 $W_c$。初始时团队策略 $\approx$ 第一阶段策略，$W_c$ 提供的修正逐步增大。
\end{remark}


% ============================================================
\section{稳定性分析}
% ============================================================

\subsection{代价函数的下界}

记 $S = \diag(s)$，组级缩放矩阵
\[
\bar{Q} = \diag(\underbrace{S^{-T}QS^{-1},\, \dots,\, S^{-T}QS^{-1}}_{n}),
\]
则
\begin{equation}
\sum_{j=1}^{n} P_j(\tilde{x}_j) = Z^T \bar{Q}\,Z \ge \lambda_{\min}(\bar{Q})\,\|Z\|^2 \triangleq \lambda_Q\,\|Z\|^2.
\end{equation}

由于 $\|u^e\|_\infty \le \bar{u}_e$，存在 $\bar{W}_e > 0$ 使 $\Psi_{\bar{u}_e, R_2}(u^e) \le \bar{W}_e$。又 $\Psi_{\bar{u}_p, R_1} \ge 0$ 且 $\gamma S_j \ge 0$，故

\begin{equation}\label{eq:cost-lower-bound}
\sum_{j=1}^{n} L_j^{\mathrm{team}} \ge \lambda_Q\,\|Z\|^2 - \bar{W}_e.
\end{equation}

\subsection{理想值函数下的 UUB}

取 Lyapunov 函数
\[
V_{\mathrm{tot}}(Z) = \sum_{j=1}^{n} V_j^{\mathrm{team},*}.
\]

\begin{assumption}\label{ass:V-pd}
$V_{\mathrm{tot}}(Z)$ 正定且径向无界。
\end{assumption}

沿理想闭环轨迹，对每个 $j$：
\[
\frac{d V_j^{\mathrm{team}}}{dt}
= \left(\frac{\partial V_j^{\mathrm{team}}}{\partial x_j^p}\right)^{\!T} \dot{x}_j^p
+ \sum_{k \neq j} \left(\frac{\partial V_j^{\mathrm{team}}}{\partial \xi_k}\right)^{\!T} \dot{\xi}_k.
\]

第一项由 HJI 鞍点条件给出 $= -L_j^{\mathrm{team}}$。

第二项为队友状态和控制变化引起的附加项。由于动力学有界（$\dot{x}_k^p$ 有界）且控制饱和（$|\dot{u}_k^p|$ 在实际中有界），存在 $M_j > 0$ 使
\begin{equation}
\left|\sum_{k \neq j} \left(\frac{\partial V_j^{\mathrm{team}}}{\partial \xi_k}\right)^{\!T} \dot{\xi}_k\right| \le M_j.
\end{equation}

记 $M_{\mathrm{tot}} = \sum_j M_j$，得
\begin{equation}\label{eq:V-dot-bound}
\dot{V}_{\mathrm{tot}} \le -\sum_{j} L_j^{\mathrm{team}} + M_{\mathrm{tot}} \le -\lambda_Q\,\|Z\|^2 + \bar{W}_e + M_{\mathrm{tot}}.
\end{equation}

\begin{theorem}[理想团队闭环 UUB]\label{thm:ideal-uub}
在假设~\ref{ass:V-pd} 下，理想团队闭环系统一致最终有界。所有轨迹最终进入
\[
\mathcal{B}_1 = \left\{Z : \|Z\| \le \sqrt{\frac{\bar{W}_e + M_{\mathrm{tot}}}{\lambda_Q}}\right\}
\]
并保持在其中。
\end{theorem}

\begin{proof}
当 $\|Z\| > \sqrt{(\bar{W}_e + M_{\mathrm{tot}})/\lambda_Q}$ 时，$\dot{V}_{\mathrm{tot}} < 0$。$V_{\mathrm{tot}}$ 递减直至轨迹进入 $\mathcal{B}_1$。
\end{proof}

\begin{remark}[UUB 而非渐近稳定的原因]
运行代价含 $-\Psi_{\bar{u}_e}(u^e)$，逃避者控制有界但一般不为零，因此全局上只能保证 UUB。只有当逃避者控制归零时才可得渐近稳定。
\end{remark}

\subsection{近似值函数下的 UUB}

实际使用近似权重 $\hat{W}_j$、$\hat{W}_c$。

\begin{assumption}\label{ass:grad-error}
梯度近似误差有界：存在 $\bar{\varepsilon}_j > 0$ 使
\[
\left\|\frac{\partial \hat{V}_j^{\mathrm{team}}}{\partial x_j^p} - \frac{\partial V_j^{\mathrm{team},*}}{\partial x_j^p}\right\| \le \bar{\varepsilon}_j, \qquad j = 1, \dots, n.
\]
\end{assumption}

由 $\tanh(\cdot)$ 的 1-Lipschitz 性，控制误差满足
\begin{equation}
\|u_j^p - u_j^{p*}\| \le c_p\,\bar{\varepsilon}_j, \qquad c_p = \frac{1}{2}\|R_1^{-1}\|\,\bar{g}_p,
\end{equation}
其中 $\bar{g}_p = \sup\|g_j(x_j^p)\|$。

近似控制引入等效扰动 $d_V(Z)$，$\|d_V\| \le \bar{d}_V$。利用 Young 不等式 $ab \le \frac{\eta}{2}a^2 + \frac{1}{2\eta}b^2$，取 $a = \|Z\|$，$b = L_V\bar{d}_V$：
\begin{equation}\label{eq:V-dot-approx}
\dot{V}_{\mathrm{tot}}
\le -\left(\lambda_Q - \frac{\eta}{2}\right)\|Z\|^2 + \bar{W}_e + M_{\mathrm{tot}} + \frac{L_V^2\,\bar{d}_V^2}{2\eta},
\end{equation}
其中 $L_V$ 为 $\nabla V_{\mathrm{tot}}$ 的局部线性增长常数。

\begin{theorem}[近似团队闭环 UUB]\label{thm:approx-uub}
在假设~\ref{ass:grad-error} 下，选 $0 < \eta < 2\lambda_Q$，近似团队闭环系统一致最终有界，终极界为
\[
\mathcal{B}_2 = \left\{Z : \|Z\| \le \sqrt{\frac{\bar{W}_e + M_{\mathrm{tot}} + L_V^2\bar{d}_V^2/(2\eta)}{\lambda_Q - \eta/2}}\right\}.
\]
\end{theorem}


% ============================================================
\section{通信退化与恢复分析}
% ============================================================

\subsection{本地估计与一致性误差}

\begin{definition}
第 $j$ 个追逐者维护队友 $k$ 的本地估计 $\hat{\xi}_{jk}$，估计误差为
\[
e_{jk}^\xi = \hat{\xi}_{jk} - \xi_k.
\]
\end{definition}

更新律：
\begin{equation}\label{eq:estimate-update}
\hat{\xi}_{jk}^+ =
\begin{cases}
\xi_k = (x_k^p, u_k^p), & \text{全通信}, \\
(x_k^p, \bm{0}), & \text{仅状态通信}, \\
\hat{\xi}_{jk}, & \text{完全断联}.
\end{cases}
\end{equation}

\subsection{全通信下的一步一致性}

\begin{proposition}\label{prop:full-consistency}
若在时刻 $t$ 对所有 $k, j$ 都有全通信，则更新后
$e_{jk}^{\xi,+} = 0$ 对所有 $j, k$ 成立。
\end{proposition}

\begin{proof}
由 \eqref{eq:estimate-update} 第一种情况直接得 $\hat{\xi}_{jk}^+ = \xi_k$。
\end{proof}

\begin{corollary}[全通信控制一致性]
当 $e_{jk}^\xi = 0$ 时，分布式执行控制与基于真实信息的集中式计算完全一致。
\end{corollary}

\subsection{状态通信下的稳定性}

仅收到 $x_k^p$ 时，$u_k^p$ 不可用，耦合特征退化为
\[
\psi_j^{\mathrm{state}} = \frac{1}{|\mathcal{N}_j|} \sum_{k \in \mathcal{N}_j} \begin{bmatrix} \phi_s(\Delta x_{jk}) \\ \bm{0}_3 \end{bmatrix}.
\]

与全通信的差异为
\[
\delta\psi_j = \psi_j - \psi_j^{\mathrm{state}} = \frac{1}{|\mathcal{N}_j|} \sum_{k \in \mathcal{N}_j} \begin{bmatrix} \bm{0}_6 \\ \phi_u(\Delta v_{jk}, u_k^p) \end{bmatrix}.
\]

由控制饱和 $\|u_k^p\| \le \sqrt{3}\,\bar{u}_p$ 及 $\Delta v$ 有界，$\delta\psi_j$ 有界。引入的控制扰动满足
\begin{equation}
\|d_u\| \le \|W_c^u\|\,\bar{g}_p\,\gamma_u\,\sqrt{3}\,\bar{u}_p\,C_v,
\end{equation}
其中 $C_v$ 为相对速度的上界。系统仍满足 UUB，终极界比全通信大。

\subsection{完全断联下的稳定性}

\begin{proposition}[退化为原文 V-SNAC]\label{prop:no-comm}
当 $\mathcal{N}_j = \varnothing$ 时，$\psi_j = \bm{0}$，价值函数退化为
\[
V_j^{\mathrm{team}} = W_j^T\phi(\tilde{x}_j),
\]
控制律退化为
\[
u_j^p = -\bar{u}_p\,\tanh\!\left(\frac{1}{2\bar{u}_p}\,R_1^{-1}\,g_j^T\left(\frac{\partial\phi}{\partial\tilde{x}_j}\right)^{\!T} W_j\right),
\]
即原文 V-SNAC。由假设~\ref{ass:stage1-converge}，系统仍满足 UUB。
\end{proposition}

\subsection{恢复后的重新一致性}

\begin{theorem}[有限时长断联后的恢复]\label{thm:recovery}
若通信中断时间有限，且恢复后满足全通信条件，则：
\begin{enumerate}[label=(\roman*)]
\item 恢复后第一次更新即实现 $e_{jk}^\xi = 0$；
\item 所有由估计误差引起的控制扰动在恢复后立刻消失；
\item 系统重新进入全通信 UUB 邻域。
\end{enumerate}
\end{theorem}

\begin{proof}
(i) 由命题~\ref{prop:full-consistency}。(ii) $e_{jk}^\xi = 0$ 时，$\psi_j$ 使用真实数据，扰动归零。(iii) 系统退化为近似全通信闭环，由定理~\ref{thm:approx-uub} 保证 UUB。
\end{proof}

\subsection{分层退化定理}

\begin{theorem}[三级退化的稳定性]\label{thm:hierarchical}
设近似团队闭环系统在全通信下满足 UUB，终极界为 $\mathcal{B}_1$。定义
\begin{align}
\mathcal{B}_2 &: \text{状态通信终极界}（丢失控制耦合 $\phi_u$），\\
\mathcal{B}_3 &: \text{完全断联终极界}（丢失全部耦合 $\psi_j$）。
\end{align}
则：
\begin{enumerate}[label=(\roman*)]
\item 状态通信下系统 UUB，终极界 $\mathcal{B}_2 \supseteq \mathcal{B}_1$；
\item 完全断联下系统 UUB，终极界 $\mathcal{B}_3 \supseteq \mathcal{B}_2$；
\item 恢复全通信后，系统重新收敛至 $\mathcal{B}_1$。
\end{enumerate}

终极界满足
\[
\mathcal{B}_1 \subseteq \mathcal{B}_2 \subseteq \mathcal{B}_3.
\]
\end{theorem}

\begin{proof}
每一级通信退化等价于引入有界扰动：
\begin{itemize}
\item 全通信 $\to$ 状态通信：扰动 $d_u$（丢失控制耦合），$\|d_u\|$ 有界；
\item 状态通信 $\to$ 断联：扰动 $d_s$（丢失状态耦合），$\|d_s\|$ 有界。
\end{itemize}
由 UUB 系统对有界扰动的鲁棒性（标准 Lyapunov 论证），每级扰动扩大终极界但不破坏有界性。

扰动大小关系：$\|d_u\| < \|d_u + d_s\|$（丢失更多信息 $\Rightarrow$ 扰动更大），故 $\mathcal{B}_1 \subseteq \mathcal{B}_2 \subseteq \mathcal{B}_3$。

恢复全通信后，$d_u = d_s = 0$，系统重新退化为定理~\ref{thm:approx-uub}，终极界回缩至 $\mathcal{B}_1$。
\end{proof}


% ============================================================
\section{与原文的对比总结}
% ============================================================

\begin{table}[htbp]
\centering
\caption{原文 V-SNAC 与本文团队价值函数的对比}
\label{tab:comparison}
\begin{tabular}{@{}lcc@{}}
\toprule
& \textbf{原文 V-SNAC} & \textbf{本文} \\
\midrule
价值函数 & $W_j^T\phi(\tilde{x}_j)$ & $W_j^T\phi(\tilde{x}_j) + W_c^T\psi_j$ \\
队友感知 & 无 & 状态 + 控制 \\
路径冲突避免 & 无 & 分离惩罚 + 耦合梯度 \\
通信断联 & 不适用 & 三级退化 \\
LS 维度 & 6 & 6（第一阶段）+ 9（第二阶段） \\
额外参数 & 0 & 9（共享 $W_c$） \\
控制律形式 & $-\bar{u}\tanh(\cdots)$ & $-\bar{u}\tanh(\cdots)$（相同） \\
稳定性 & UUB & UUB（含通信扰动项） \\
\bottomrule
\end{tabular}
\end{table}


% ============================================================
\appendix
\section{追逐者控制律推导细节}\label{app:control}
% ============================================================

对第 $j$ 个追逐者，Hamiltonian 中与 $u_j^p$ 相关的驻值条件为
\[
\frac{\partial H_j}{\partial u_j^p} = \frac{\partial \Psi_{\bar{u}_p, R_1}}{\partial u_j^p} + g_j(x_j^p)^T \frac{\partial V_j^{\mathrm{team}}}{\partial \tilde{x}_j} = 0.
\]

对非二次代价 \eqref{eq:nonquad} 的单分量
\[
\varphi(u) = 2\bar{u}\,r\left[u\,\atanh\!\left(\frac{u}{\bar{u}}\right) + \frac{\bar{u}}{2}\ln\!\left(1 - \frac{u^2}{\bar{u}^2}\right)\right],
\]
求导得
\[
\frac{d\varphi}{du} = 2\bar{u}\,r\,\atanh\!\left(\frac{u}{\bar{u}}\right).
\]

因此
\[
\frac{\partial\Psi_{\bar{u}_p, R_1}}{\partial u_j^p} = 2\bar{u}_p\,R_1\,\atanh\!\left(\frac{u_j^p}{\bar{u}_p}\right).
\]

代入驻值条件：
\[
2\bar{u}_p\,R_1\,\atanh\!\left(\frac{u_j^{p*}}{\bar{u}_p}\right) + g_j^T \frac{\partial V_j^{\mathrm{team}}}{\partial \tilde{x}_j} = 0.
\]

整理：
\[
\atanh\!\left(\frac{u_j^{p*}}{\bar{u}_p}\right) = -\frac{1}{2\bar{u}_p}\,R_1^{-1}\,g_j^T\,\frac{\partial V_j^{\mathrm{team}}}{\partial \tilde{x}_j}.
\]

两边取 $\tanh$（利用 $\tanh(-x) = -\tanh(x)$）：
\[
u_j^{p*} = -\bar{u}_p\,\tanh\!\left(\frac{1}{2\bar{u}_p}\,R_1^{-1}\,g_j^T\,\frac{\partial V_j^{\mathrm{team}}}{\partial \tilde{x}_j}\right).
\]

由 $\partial\tilde{x}_j/\partial x_j^p = I_6$，得 $\partial V_j^{\mathrm{team}}/\partial\tilde{x}_j = \partial V_j^{\mathrm{team}}/\partial x_j^p$，此即 \eqref{eq:team-up}。\qed

\bigskip

\paragraph{逃避者控制律推导}

对逃避者，Hamiltonian 的总和中与 $u^e$ 相关的部分：
\[
\frac{\partial}{\partial u^e}\sum_{j=1}^{n} H_j = -\frac{\partial\Psi_{\bar{u}_e, R_2}}{\partial u^e} + \sum_{j=1}^{n}\left(\frac{\partial V_j^{\mathrm{team}}}{\partial\tilde{x}_j}\right)^{\!T}(-g_e) = 0.
\]

整理：
\[
2\bar{u}_e\,R_2\,\atanh\!\left(\frac{u^{e*}}{\bar{u}_e}\right) + g_e^T \sum_{j=1}^{n}\frac{\partial V_j^{\mathrm{team}}}{\partial\tilde{x}_j} = 0.
\]

同理取 $\tanh$ 得
\[
u^{e*} = -\bar{u}_e\,\tanh\!\left(\frac{1}{2\bar{u}_e}\,R_2^{-1}\,g_e^T \sum_{j=1}^{n}\frac{\partial V_j^{\mathrm{team}}}{\partial\tilde{x}_j}\right).
\]
\qed


\end{document}
